\section{综合结果分析}

表~\ref{tab:overall} 汇总各问的正式结果及其对照。问题一为典型日，问题二至四为 $2025$ 年 $2$--$12$ 月共 $334$ 天；费用降幅均相对于本问同口径基准计算。

\begin{table}[!htbp]
  \centering
  \caption{四问关键结果对照}
  \label{tab:overall}
  \small
  \setlength{\tabcolsep}{4pt}
  \begin{tabularx}{\textwidth}{@{}l>{\raggedright\arraybackslash}X>{\raggedright\arraybackslash}X>{\raggedright\arraybackslash}X>{\raggedright\arraybackslash}X@{}}
    \toprule
    问题 & 适用条件与新增机制 & 本文方案 & 关键费用（元） & 同口径对照与节省 \\
    \midrule
    一 & 典型日、确定性日前计划 & 线性规划 & 日购电费 $35\,126.85$ & 无储能 $48\,052.05$，省 $26.90\%$ \\
    \addlinespace
    二 & 全年日前计划；负载/光伏预测；$5$ 倍紧急购电 & 净分位 $+$ 日末 $2400$ $+$ $\rho$ 回放 & 全年 $13\,695\,403.61$ & 固定 $0.8$ 为 $13\,765\,167.58$，省 $0.51\%$ \\
    \addlinespace
    三 & 日内滚动调整；少买 $0.5$ 倍、多买 $1.5$ 倍偏差结算 & $48$ h 展望 $+$ 分源残差修正 & 全年 $13\,275\,795.85$ & 不调整 N0 $15\,197\,639.94$，省 $12.6\%$ \\
    \addlinespace
    四 & 实时波动电价；决策用 $\hat p_0$、结算用附件 4 & 三层电价预测 $+$ 净分位 $+\rho$ / LA $q{=}0.60$ & Q4-2 $14\,413\,479.02$；Q4-3 $14\,001\,199.26$ & 固定 $14\,480\,739.96$，省 $0.46\%$；N0 $15\,983\,608.51$，省 $12.4\%$ \\
    \bottomrule
  \end{tabularx}
\end{table}

由表~\ref{tab:overall} 可见，线性规划在确定性信息下统筹全天调度；预测条件下，日内滚动调整较不调整降低了费用，这一方向在两种电价情景下均成立。各问的信息条件与电价不同，费用绝对值不直接用于比较模型优劣。

参数标定和可行性检验见各问。功率口径对照、旧方案消融及完美信息诊断集中保存在支撑材料 \texttt{condensed\_sections\_notes.md}，其中历史费用不替代本表的正式结果。
